\subsection{Additional results}
\label{app:additional-results}

This appendix contains figures that support the discussion in
\cref{sec:results-discussion}, but are not needed to follow it.

\begin{figure}[htbp]
     \centering
     \includegraphics[width=\linewidth]{figures/regression/history.pdf}
     \caption{Training MSE per epoch for the selected FFNNs: one hidden layer of 50 nodes
     (Adam, $\gamma=0.1$, 100 epochs) and two hidden layers of 100 nodes (RMSprop,
     $\gamma=10^{-4}$, 10000 epochs), both with a batch size of 32. The larger network
     stays at the mean prediction, with an MSE near the variance of the targets, for a large
     part of the training.}
     \label{fig:training-history}
\end{figure}

\begin{figure}[htbp]
     \centering
     \includegraphics[width=\linewidth]{figures/regression/activation_comparison.pdf}
     \caption{Predictions of the best sigmoid, ReLU and Leaky ReLU networks from
          \cref{tab:activation-comparison}, together with the Runge function and the
          noisy training and test data. The ReLU and Leaky ReLU curves are nearly
          identical and consist of a few straight segments, while the sigmoid curve
          is smooth.}
     \label{fig:activation-comparison}
\end{figure}

\begin{figure}[htbp]
     \centering
     \includegraphics[width=\linewidth]{figures/regression/regularization_comparison.pdf}
     \caption{The best $\ell_2$ and $\ell_1$ regularized networks, compared with the
          Ridge and Lasso models of Project~1, together with the Runge function and the
          noisy training and test data. Ridge oscillates near $x = 1$ and Lasso underfits
          the peak, while both networks follow the function closely.}
     \label{fig:regularization-comparison}
\end{figure}

\begin{figure}[htbp]
     \centering
     \includegraphics[width=\linewidth]{figures/classification/logreg_confusion_matrix.pdf}
     \caption{Confusion matrix of logistic regression on the test set, normalized by
     row so that the diagonal is the recall of each digit (\%). Most errors are between
     digits that share strokes, such as 3, 5 and 8.}
     \label{fig:logreg-matrix}
\end{figure}

\begin{figure}[htbp]
     \centering
     \includegraphics[width=\linewidth]{figures/classification/ffnn_regularization.pdf}
     \caption{Validation accuracy (\%) of the selected architecture for each learning rate
     and $\ell_2$ penalty. A small penalty helps slightly, while a large one, especially
     with a large learning rate, makes training fail.}
     \label{fig:ffnn-regularization}
\end{figure}

\begin{figure*}[htbp]
     \centering
     \includegraphics[width=0.8\textwidth]{figures/regression/regularization_heatmaps.pdf}
     \caption{Cross-validated MSE over the penalty $\lambda$ and the learning rate $\gamma$,
          for the architecture of the best configuration of each activation and norm. The
          framed cell is the minimum. The ReLU family has its minimum at $\lambda = 10^{-3}$
          for both norms, while every sigmoid network collapses to predicting the mean for
          all but the smallest penalties.}
     \label{fig:regularization-heatmaps}
\end{figure*}
